Под \emph{симплициальным комплексом} мы будем понимать конечный набор $K$ симплексов, замкнутый относительно взятия граней: если $\sigma \in K$ и $\tau \subseteq \sigma$, то $\tau \in K$. Симплициальный комплекс служит дискретным аналогом топологического пространства и базовым объектом, на котором будет считаться гомология.

Когда данные приходят в виде облака точек в метрическом пространстве, чтобы превратить их в симплициальный комплекс, обычно используют конструкцию комплекса Чеха.
\begin{definition}[Комплекс Чеха]
Пусть $(X, d)$ - метрическое пространство, $P \subseteq X$ - конечное множество, и $\varepsilon \ge 0$.
\emph{Комплексом Чеха} $\check{C}_{\varepsilon}(P)$ называется симплициальный комплекс с множеством вершин $P$, в котором
конечное подмножество $\sigma\subseteq P$ является симплексом тогда и только тогда, когда
\[
\bigcap_{p\in \sigma} B(p, \varepsilon) \neq \varnothing,
\]
где $B(p, \varepsilon)$ - открытый шар радиуса $\varepsilon$ с центром в точке $p$.
\end{definition}

Реальные данные обычно зашумлены, поэтому один комплекс Чеха, построенный по фиксированному масштабу $\varepsilon$, говорит о геометрии облака немного. Кроме того, каждая точка может нести дополнительную информацию (например, значение какой-то функции), которая в одиночном комплексе Чеха никак не учитывается. Чтобы преодолеть оба ограничения, вводится понятие фильтрации --- семейства комплексов, упорядоченных по шкале параметра.
\begin{definition}[Фильтрация]
Пусть $K$ --- симплициальный комплекс, а $(I,\le)$ --- частично упорядоченное множество.
\emph{Фильтрацией} комплекса $K$, индексированной $I$, называется семейство симплициальных подкомплексов
$\{K_\alpha\}_{\alpha\in I}$, где $K_\alpha\subseteq K$, такое что для любых $\alpha\le \beta$ выполнено
\[
K_\alpha \subseteq K_\beta.
\]
Мы будем брать $I=\mathbb{R}$ или $I=\mathbb{Z}$ с обычным порядком.
\end{definition}

\begin{examples}
Приведем важнейшие примеры фильтраций
\begin{itemize}
\item Семейство комплексов Чеха $\{\check{C}_{\varepsilon}(P)\}_{\varepsilon \ge 0}$, параметризованное $\varepsilon \in \mathbb{R}_{\ge 0}$, является фильтрацией: из $\varepsilon \le \varepsilon'$ следует $\check{C}_{\varepsilon}(P)\subseteq \check{C}_{\varepsilon'}(P)$.
Такой объект даёт значительно более репрезентативное представление облака точек, чем комплекс Чеха с фиксированным масштабом.

\item Фильтрация по подуровням функции.
Пусть $K$ - симплициальный комплекс с множеством вершин $V$, и пусть задана функция $w:V\to\mathbb{R}$.
Продлим её на симплексы по правилу
\[
w(\sigma) := \max_{v\in\sigma} w(v).
\]
Для $a\in\mathbb{R}$ положим
\[
K_a := \{\sigma\in K : w(\sigma)\le a\}.
\]
Тогда каждое $K_a$ является симплициальным подкомплексом, и $\{K_a\}_{a\in\mathbb{R}}$ образует фильтрацию комплекса $K$. Это позволяет учитывать скрытое представление каждой точки. Например, изображение как облако пикселей в плоскости само по себе малоинформативно, но если задать функцию из каждого пикселя в его яркость и построить по ней фильтрацию по подуровням, мы получим осмысленное топологическое описание изображения.
\end{itemize}
\end{examples}

% \begin{definition}[Симплициальные гомологии]
% Пусть $K$ - симплициальный комплекс. Для $k\ge 0$ определим $\Bbbk$-векторное пространство $k$-цепей:
% \[
% C_k(K;\Bbbk) := \bigoplus_{\sigma\in K_k} \Bbbk\cdot \sigma,
% \]
% то есть свободное $\Bbbk$-пространство, порождённое (ориентированными) $k$-симплексами.
% Граничный оператор $\partial_k: C_k(K;\Bbbk)\to C_{k-1}(K;\Bbbk)$ задаётся на порождающих формулой
% \[
% \partial_k [v_0,\dots,v_k] := \sum_{i=0}^k (-1)^i [v_0,\dots,\widehat{v_i},\dots,v_k],
% \]
% и продолжается $\Bbbk$-линейно. Выполнено $\partial_{k-1}\circ\partial_k=0$.
% Определим пространства циклов и границ:
% \[
% Z_k(K;\Bbbk):=\ker(\partial_k),\qquad B_k(K;\Bbbk):=\operatorname{im}(\partial_{k+1}).
% \]
% Тогда $k$-я симплициальная группа гомологий равна
% \[
% H_k(K;\Bbbk) := Z_k(K;\Bbbk)/B_k(K;\Bbbk).
% \]
% \end{definition}

\begin{definition}[Персистентные гомологии]
Пусть $\{K_a\}_{a \in I}$ --- фильтрация, индексированная линейно упорядоченным множеством $(I, \le)$, и пусть $n \ge 0$.
Для $a\le b$ включение $K_a\hookrightarrow K_b$ индуцирует гомоморфизм цепных комплексов и, соответственно, линейное отображение гомологий:
\[
f^{(n)}_{a\le b}: H_n(K_a;\Bbbk)\to H_n(K_b;\Bbbk).
\]
\emph{$n$-й персистентной гомологией} от $a$ к $b$ называется образ
\[
H_n^{a\to b}(\{K_\alpha\}_{\alpha \in I};\Bbbk) := \operatorname{Im}\bigl(f^{(n)}_{a\le b}\bigr)\subseteq H_n(K_b;\Bbbk).
\]
Эквивалентно,
\[
H_n^{a\to b} \cong \frac{Z_n(K_a;\Bbbk)}{Z_n(K_a;\Bbbk)\cap B_n(K_b;\Bbbk)}.
\]
Система $\{H_n(K_a;\Bbbk), f^{(n)}_{a\le b}\}_{a\le b \in I}$ называется $n$-м персистентным модулем фильтрации $\{K_a\}_{a \in I}$.
\end{definition}

\begin{definition}[Диаграмма персистентности]
В случае $I = \{0, 1, \dots, N\}$ прямую сумму $\bigoplus_{a \in I} H_n(K_a, \Bbbk)$ можно рассматривать как модуль над кольцом $\Bbbk[t]$, где действие $t$ задаётся правилом
\[
v \in H_n(K_a; \Bbbk), \quad t \cdot v = f^{(n)}_{a \le a+1}(v).
\]
Тогда по структурной теореме о конечно порождённых модулях над областями главных идеалов он раскладывается в прямую сумму интервальных модулей вида $[b_i, d_i)$, где $b_i < d_i \le \infty$.

\emph{$n$-й диаграммой персистентности} $\mathrm{Dgm}_n$ называется мультимножество точек
\[
\mathrm{Dgm}_n := \{(b_i, d_i)\}_i \subseteq \overline{\mathbb{R}}^2, \qquad \overline{\mathbb{R}} := \mathbb{R} \cup \{\infty\},
\]
взятое с кратностями, соответствующее интервалам в указанном разложении. Точки вблизи диагонали $b = d$ принято считать шумом, тогда как точки с большой персистентностью $d - b$ отвечают устойчивым топологическим признакам.
\end{definition}
